\documentclass[11pt,a4paper]{article}
\usepackage[margin=1in]{geometry}
\usepackage[T1]{fontenc}
\usepackage[utf8]{inputenc}
\usepackage{lmodern}
\usepackage{microtype}
\usepackage{hyperref}
\usepackage{xurl}
\usepackage{enumitem}
\usepackage{amsmath}

\setlength{\emergencystretch}{2em}

\title{Design Report: Satellite Debris Collision-Avoidance Runtime and RISC-V Emulator}
\author{HackRush 2026 -- Computer Architecture Project}
\date{April 2026}

\begin{document}
\maketitle

\section{Introduction}
This report summarizes the design of a deterministic satellite debris collision-avoidance runtime and its companion RISC-V emulator. The primary engineering goal is to make scheduling, risk decisions, cycle accounting, and energy accounting transparent and reproducible. Instead of optimizing only for realism, the implementation prioritizes traceability: every major runtime decision has an explicit formula, fixed interval, or bounded state transition that can be audited directly in logs.

The implementation has two complementary parts. First, a C11 runtime executes periodic mission tasks (orbit propagation, risk evaluation, maneuver planning, telemetry, and transmission) under a fixed cycle budget per second. Second, a Python emulator executes RV32-style assembly kernels with pipeline-aware cycle accounting and lightweight power estimation. Together, they support both system-level reasoning (mission behavior, deadlines, energy budget compliance) and architecture-level reasoning (instruction mix, hazards, CPI behavior).

Key design intent for this submission is practical and specific: use integer/fixed-point math for deterministic behavior, use simple but meaningful risk thresholds tied to debris size, use small bounded maneuver candidate sets for predictable compute cost, and expose enough micro-architectural detail in the emulator to validate cycle/energy trends against assembly scenarios.

\section{Emulator Design and Hardware Model Assumptions}
The emulator in \texttt{simulator\_core/emulator/emulator.py} uses the following model:
\begin{itemize}[leftmargin=1.5em]
    \item RV32-style integer register file (32 registers with ABI aliases).
    \item Instruction subset needed by project assembly kernels:
    \texttt{li, mv, addi, add, sub, ori, mul, div, rem, slt, lw, sw, beq, bne, blt, bge, bnez, beqz, j, jal, ret}.
    \item Bounded memory model with text/data/stack regions and optional fast-memory abstraction.
\end{itemize}
This model is intentionally scoped to the project kernels: it is detailed enough for instruction-level timing discussion while avoiding unrelated ISA complexity.

Timing assumptions:
\begin{equation}
\texttt{total\_cycles} = (k-1) + \sum_i \texttt{cost}(i)
\end{equation}
where $k=5$ (pipeline fill overhead), and \texttt{cost(i)} includes base latency plus penalties.

Default baseline currently used:
\begin{itemize}[leftmargin=1.5em]
    \item ALU/load/store/branch/mul/div/rem base cost = 1,
    \item taken-branch penalty = 0,
    \item memory penalty = 0,
    \item load-use hazard stall = +1 when detected.
\end{itemize}
These defaults produce near-ideal throughput but still capture meaningful hazards. In practice, the load-use check is the most relevant penalty for the provided kernels, because dependent instruction pairs after \texttt{lw} appear naturally in state-update code.

Power model in emulator:
\begin{equation}
E_{mJ} = C_{active}P_{active}T + C_{sleep}P_{sleep}T + C_{radio}P_{radio}T
\end{equation}
with configurable defaults in code.
The reported CPI is derived from total cycles, so it remains consistent with the pipeline-fill model ($n + k - 1$) plus any dynamic penalties.

\section{OS/Runtime Design}
\subsection*{Scheduler type}
The runtime in \texttt{simulator\_core/src/main.c} uses a static time-triggered cyclic scheduler over 1-second ticks:
\begin{itemize}[leftmargin=1.5em]
    \item Every 5 s: orbit propagation.
    \item Every 10 s: risk evaluation and maneuver.
    \item Every 20 s: telemetry and transmission.
\end{itemize}
Release checks are modulo-based and constant-time. This makes dispatch deterministic and easy to validate through timestamped logs.

\subsection*{Priorities, timing, and deadlines}
No dynamic priority queue is used; effective priority is release order in each tick:
\begin{enumerate}[leftmargin=1.5em]
    \item Orbit propagation,
    \item Risk evaluation,
    \item Maneuver,
    \item Telemetry,
    \item Transmission.
\end{enumerate}
Each second has a fixed cycle budget (\texttt{SIM\_CYCLES\_PER\_SEC=100000}).
Used cycles are charged to tasks; leftover cycles are charged to sleep.
This yields deterministic accounting and bounded per-tick utilization.

Task compute costs are explicitly modeled in cycles and are central to timing analysis:
\begin{itemize}[leftmargin=1.5em]
    \item Orbit propagation: \(1800 + 25 \times \texttt{debris\_count}\).
    \item Risk evaluation: \(4200 + 70 \times \texttt{debris\_count}\).
    \item Maneuver: 2200 when a maneuver is pending, else 300.
    \item Telemetry: 1100.
    \item Transmission: 800.
\end{itemize}
Because these costs are deterministic functions of current state (especially \texttt{debris\_count}), timing behavior can be reasoned about analytically and then verified directly against output logs.

\subsection*{Power-management policy}
Implemented in \texttt{simulator\_core/src/runtime.c}:
\begin{itemize}[leftmargin=1.5em]
    \item Active task cycles add active energy.
    \item Radio-on tasks add extra radio energy.
    \item Remaining per-tick capacity is counted as sleep (low-power state).
\end{itemize}
The energy constants used by the runtime are:
\begin{itemize}[leftmargin=1.5em]
    \item Active work: 200 units/cycle \(=0.020\,\text{mJ/cycle}\).
    \item Sleep: 5 units/cycle \(=0.0005\,\text{mJ/cycle}\).
    \item Radio extra: 500 units/cycle \(=0.050\,\text{mJ/cycle}\).
\end{itemize}
One internal unit is \(10^{-4}\) mJ (tracked as \texttt{energy\_mj\_x10000}).
This fixed-point energy representation avoids floating-point drift and keeps budget checks stable across builds and platforms.

\section{Architecture-Aware Optimizations}
\subsection*{Fixed-point arithmetic: where and why}
Fixed-point/integer math is used to reduce overhead and improve deterministic behavior:
\begin{itemize}[leftmargin=1.5em]
    \item Angles: millidegrees (\texttt{mdeg}).
    \item Energy: \texttt{energy\_mj\_x10000} (0.0001 mJ units).
    \item Integer rad-to-mdeg conversion constant (57296).
\end{itemize}
This choice is critical for two reasons: frequent execution in the scheduler loop, and the need for deterministic comparisons in risk classification and budget accounting. Integer/fixed-point arithmetic keeps both predictable.

\subsection*{Assembly: where and why}
Assembly in \texttt{simulator\_core/asm} is used for architecture validation and instruction-level analysis:
\begin{itemize}[leftmargin=1.5em]
    \item \texttt{runtime\_kernels\_riscv32.s}: helper kernels mirroring runtime logic.
    \item \texttt{sort\_10\_registers\_riscv32.s}: register-only sorting micro-benchmark.
\end{itemize}
These kernels are used to verify both correctness and timing metrics (instruction count, total cycles, CPI, stall behavior) under a controlled instruction stream.

\subsection*{Data layout and locality}
\begin{itemize}[leftmargin=1.5em]
    \item Compact arrays/structs for task and debris data to improve sequential access.
    \item Linear iteration over debris catalog in orbit/risk tasks.
    \item Fast-memory region parameter in emulator (scratchpad-like abstraction).
\end{itemize}
The scenario structure supports up to \texttt{MAX\_DEBRIS = 256}. Critical loops traverse the debris array linearly, which is locality-friendly in the runtime and simple to reason about in cycle-level emulator experiments.

Design-specific risk thresholds are also intentionally simple and integer-based. For each debris item:
\begin{equation}
\texttt{proximity} = |\Delta r| + \frac{|\Delta \theta|}{200}
\end{equation}
\begin{itemize}[leftmargin=1.5em]
    \item HIGH-RISK if \(\texttt{proximity} < 100 \times \texttt{size\_cm}\)
    \item WATCH if \(\texttt{proximity} < 1000 \times \texttt{size\_cm}\)
    \item SAFE otherwise
\end{itemize}

This thresholding strategy directly scales risk tolerance by debris size, which is a practical mission heuristic: larger objects are classified as risky at larger separations.

Maneuver selection evaluates a small fixed candidate set
\(\Delta v\in\{-3,+3\}\,\text{cm/s}\),
\(\Delta r\in\{-1,+1\}\,\text{m}\),
and picks the option with best predicted next-step proximity.
The bounded 2x2 search keeps compute cost constant while still allowing both tangential and radial correction.

\section{Implemented Extras Relevant to This Submission}
\begin{itemize}[leftmargin=1.5em]
    \item GUI fallback to a valid entry label when selected label is missing.
    \item Simplified emulator performance summary: instruction count, total cycles, CPI, stall cycles, and estimated energy.
    \item Added over-budget test scenario \texttt{energy\_budget\_exceeded\_scenario.json} with validated \texttt{battery\_status=EXCEEDED}.
    \item Updated project scripts and docs for renamed core folder \texttt{simulator\_core}.
\end{itemize}
These additions improve reliability of day-to-day usage (GUI loading behavior), readability of performance reporting, and reproducibility of stress testing for energy-budget violations.

\section{Conclusion}
The final design is intentionally balanced: simple enough to remain analyzable, but detailed enough to support meaningful architecture and runtime conclusions. The runtime delivers deterministic scheduling and transparent energy/cycle accounting, while the emulator provides instruction-level visibility under an explicit timing model. Most importantly, the selected thresholds, task-cost formulas, and fixed-point representations are tightly aligned with the project goal of reproducible mission-level decision making under constrained resources.

The main limitations are also clear in the current implementation: scheduling is non-preemptive at 1-second granularity, scenario parsing is lightweight rather than fully standards-compliant JSON parsing, emulator ISA coverage is limited to the project-required subset, and cache/memory hierarchy effects are abstract rather than cycle-accurate. The coarse 1-second base tick also means interrupt jitter and asynchronous event behavior are not modeled at RTOS fidelity.

Future improvements follow directly from these limitations: add preemptive scheduling with deadline-miss diagnostics, upgrade to strict JSON parsing with robust validation, extend emulator ISA and branch/memory hierarchy realism, and add automated regression checks for CPI/energy invariants and scenario-level correctness.

\end{document}
